\documentclass{article}
\usepackage{graphicx}
\usepackage{amsmath}
\usepackage{amssymb}
\usepackage{amsfonts}
\usepackage{hyperref}
\usepackage{exercises}
\usepackage{indentfirst}
\usepackage{circledsteps}
\usepackage{amsthm}
\usepackage{csquotes}
\usepackage{yhmath}
\usepackage{comment}
\usepackage{placeins}
\usepackage{geometry}
\usepackage{hyperref}
\usepackage{tikz}
\usepackage{tabto}
\newcommand*{\oldneg}{\mathord{\sim}}
\newcommand{\Z}{\mathbb{Z}}
\newcommand{\N}{\mathbb{N}}
\let\oldemptyset\emptyset

\theoremstyle{definition}
\newtheorem{definition}{Definition}
\newtheorem{theorem}{Theorem}[section]
\newtheorem{corollary}{Corollary}[theorem]
\newtheorem{lemma}[theorem]{Lemma}

\setlength{\parindent}{.5in}
\geometry{margin=0.5in}
\setlength{\parskip}{12pt}
\title{Chapter 13 Notes}
\begin{document}
\date{}
\maketitle

\section*{Section 1}

\subsection*{Notes}
Note that, on page 170, the separation axioms are listed incorrectly. They should use the 
formula $\phi(y,y_1,...,y_n)$, not $\phi(x,y_1,...,y_n)$.

\subsection*{Exercise 1.1 (a)}
We prove that $K$ has property $F'_2$ iff all the separation axioms are true over $K$.

(a) On the one hand, suppose that $K$ has property $F'_2$. Now consider any set $x\in K$ and any 
formula $\phi(y,y_1,...,y_n)$ (for some elements $y_1,...,y_n$ of $K$). Then let $A$ be the class of all $y$ such that 
$\phi(y,y_1,...,y_n)$ is true over $K$ ($A$ is defined by $\phi(y,y_1,...,y_n)$ over $K$). By 
$F'_2$, $A\cap x\in K$. But $A\cap x$ satisfies the formula in $\textbf{Ax\, 2}$. Thus, for each 
formula $\phi(y,y_1,...,y_n)$ and $x\in K$, the separation axiom for that formula and element 
holds over $K$. 

(b) On the other hand, suppose that all the separation axioms are true over $K$. Then consider any 
first-order definable subclass $A$ of $K$ and any element $b$ of $K$. $A$ must be definable over $K$ by 
some formula $\phi(x)$. Then, by the separation axiom, there is some $z\in K$ such that, for all 
$x\in K$, $x\in z\equiv(x\in b\wedge\phi(x))$. This formula works as, if 
$x\in b\cap A$, then $x\in A$, and thus $x\in K$, meaning the formula would hold true. Of course, if 
the formula is true, then $x\in b\cap A$.

\subsection*{Exercise 1.1 (b)}
We prove that $K$ (where $K$ is transitive) is a first-order universe iff all the axioms of $ZF$ are true over $K$.

(a) On the one hand, suppose that $K$ is a first-order universe. We show that all the $ZF$ axioms 
are true over $K$.
 
$\textbf{Ax\, 1}$: Suppose that $(\forall z)(z\in x\equiv z\in y)$. Then, because $K$ is transitive
(by $F_1$), it follows that $x=y$.

$\textbf{Ax\, 2}$: We show that this holds by showing that $F'_2$ holds (in $K$), and consequently 
$\textbf{Ax\, 2}$ would hold by Exercise 1.1 (a). So suppose $A$ is a first-order definable subclass of 
$K$ and $b$ is an element of $K$. Then, because $K$ is first-order swelled, it follows that 
$A\cap b\in K$ (because $A\cap b$ is first-order definable over $K$, as it is defined by the 
formula $(y\in z\equiv(y\in x\wedge\phi(y)))$ where $\phi(y)$ is some formula that defines 
$A$ over $K$).

$\textbf{Ax\, 3}$: This follows by $F_3$.

$\textbf{Ax\, 4}$: This follows by $F_4$.

$\textbf{Ax\, 5}$: This follows by $F_5$, as $K$ is transitive. This is because, as $K$ is transitive, 
the set $y$ defined by
\[ (\exists y)(\forall z)(z\in y\equiv(\exists w)(z\in w\wedge w\in x)) \] is 
precisely the union of $x$.

$\textbf{Ax\, 6}$: This follows by $F_6$. For, if $x\in K$, then 
$\mathcal{P}(x)\cap K\in K$. But $\mathcal{P}(x)\cap K$ is precisely the $y$ defined by the formula
\[ (\exists y)(\forall z)(z\in y\equiv z\subseteq x) \]

$\textbf{Ax\, 7}$: This follows by $F_7$.

$\textbf{Ax\, 8}$: This follows by $F_8$.

$\textbf{Ax\, 9}$: This follows by $F_9$.

(b) On the other hand, suppose that all the $ZF$ axioms are true over $K$. We show that 
$K$ is a first order universe.

$F_1$: This holds by assumption.

$F_2$: Consider any $x\in K$ and some first-order definable (over $K$) subset $A$ of $x$. By Lemma 1.1, 
condition $F'_2$ holds, which means that $A\cap x=A$ is in $K$. Thus $K$ is first-order swelled.

$F_3$: This follows by $\textbf{Ax\, 3}$ (because $K$ is transitive).

$F_4$: This follows by $\textbf{Ax\, 4}$ (again because $K$ is transitive).

$F_5$: This follows by $\textbf{Ax\, 5}$ (again because $K$ is transitive).

$F_6$: This follows by $\textbf{Ax\, 6}$.

$F_7$: This follows by $\textbf{Ax\, 7}$.

$F_8$: If $F$ is a function that is first-order definable over $K$, then it is defined by 
some formula $\phi(x,y)$. Then, for every $w\in K$, by $\textbf{Ax\; 8}$ there is a $v\in K$ such that
\[ (\forall y)(y\in v\equiv(\exists x)(x\in w)\wedge\phi(x,y)) \]
which is just saying $F''(w)\in K$.

$F_9$: This follows by $\textbf{Ax\, 9}$ (because $K$ is transitive).

\subsection*{Exercise 1.2}
It is not clear what is being asked here. Are the $ZF$ axioms being discussed the first-order ones, or the ones introduced at the beginning 
of the book? If they are the former, clearly we cannot be using Theorems from chapters 7 and 8, where the first-order axioms 
were not yet introduced. For the problem, I am therefore assuming we are speaking of the intial $ZF$ axioms. Well, suppose we 
add the axiom of well foundedness. By Theorem 3.5 of Chapter 8, it follows that $V=R_\Omega$, which is a Zermelo-Fraenkel Universe
by Theorem 4.3 of Chapter 7. Thus these new axioms are consistent by assumption.

This means that the axiom of well foundedness cannot be disproved in $ZF$ for the following reason: For suppose it could be. 
Then adding the axiom of well foundedness to the $ZF$ axioms would introduce a contradiction, contrary to the result of this 
exercise.

\end{document}

